\section{美国数据中心税收优惠：政策格局与近期变化}

美国的数据中心税收优惠主要由州和地方政府设立。对于加密矿企改造或新建数据中心的项目，州级税收优惠主要影响两个方面。

第一，建筑材料与设备采购。在数据中心改造过程中，矿企可能需要购入大量有形动产，包括服务器、电气设备和网络设备，以支持新的业务活动。这些采购可能触发销售和使用税，而部分州会在满足特定条件时提供免税待遇。

第二，数据中心的持续运营。矿企的数据中心竣工并投入使用后，土地、建筑和设备仍可能需要缴纳财产税，大量用电也可能涉及销售税、公用事业税或用电税。因此，部分州和地方政府会提供财产税减免、与电力相关的免税待遇或其他经济激励。

多年来，美国许多州通过税收优惠吸引数据中心投资。但到2026年，公众对数据中心基础设施的态度开始由支持转向反对，政策环境也出现明显变化。

得克萨斯州、弗吉尼亚州和佐治亚州等主要辖区的销售和使用税优惠，仍适用于符合条件的项目。与此同时，AI 数据中心快速扩张引发了对电价、水资源和当地居民负担的广泛关注，促使多个州暂停数据中心优惠，或考虑采取类似措施。自2026年夏季起，亚利桑那州对新的计算机数据中心税收减免申请实施为期三年的暂停受理措施；华盛顿州取消了服务器更换设备和翻新项目的销售和使用税优惠；路易斯安那州则提高了数据中心优惠的附加条件，包括要求数据中心企业承担其项目带来的全部新增电力需求。

这意味着，加密矿企在转型过程中仍可能获得税收优惠，但优惠的可获得性正变得更不确定。项目选址、是否需要新的免税资格认证、申请时点以及政策调整引入的新因素，都可能影响转型的实际税务成本。

% 依据英文版表1逐项翻译，数值与条件沿用原稿。
\begingroup
\ReportTableSetup
\begin{longtable}{|L{0.15\linewidth}|L{0.23\linewidth}|L{\dimexpr0.62\linewidth-6\tabcolsep-4\arrayrulewidth\relax}|}
\caption{各州数据中心税收优惠状态}\label{tab:state-incentives}\\
\hline
\textbf{州} & \textbf{当前状态} & \textbf{主要变化} \\ \hline
\endfirsthead
\caption[]{各州数据中心税收优惠状态（续）}\\
\hline
\textbf{州} & \textbf{当前状态} & \textbf{主要变化} \\ \hline
\endhead
亚利桑那州 & 暂停新申请 & 自2026年7月1日至2029年6月30日，亚利桑那州商务管理局不得受理新计算机数据中心的申请。 \\ \hline
俄亥俄州 & 暂停新申请 & 俄亥俄州税收抵免管理局已暂停审议新的数据中心税收豁免请求。 \\ \hline
伊利诺伊州 & 暂停新协议 & 自2026年7月1日起，该州停止处理数据中心激励计划的新申请。 \\ \hline
马萨诸塞州 & 暂停新申请 & 自2026年6月25日起，该州暂停受理合资格数据中心销售和使用税豁免申请。 \\ \hline
华盛顿州 & 缩小免税范围 & 自2026年7月1日起，服务器更换设备不再适用原有豁免；不再签发新的农村地区翻新证书，城市地区翻新项目也不再符合资格。 \\ \hline
弗吉尼亚州 & 仍可免税；新增用电税 & 销售和使用税豁免继续有效；自2026年7月1日至2028年6月30日，按每千瓦时0.011美元征收数据中心用电税。 \\ \hline
路易斯安那州 & 仍可免税；条件待定 & 2026年6月的一项行政命令要求，为数据中心销售税优惠建立电费缴纳者和社区保护框架，具体标准由相关主管部门制定中。 \\ \hline
得克萨斯州 & 仍可免税；政策变化风险较高 & 该州要求全面审查数据中心电网互联项目，并提出与立法机关合作，取消过时或不必要的数据中心优惠。 \\ \hline

\end{longtable}
\endgroup


\FloatBarrier
